\section{Results}
\label{sec:results}

\subsection{Predictive Performance and Decision Bands}
Table~\ref{tab:models} compares the final model with the baselines. The final LightGBM model reaches a test ROC-AUC
of 0.767 and a PR-AUC of 0.255, against 0.748 for the best baseline (XGBoost trained on SMOTE-balanced data). Excluding
gender lowers the ROC-AUC by only 0.0015. The ten bootstrap-retrained models reach $0.7609 \pm 0.0009$; they are
slightly weaker because each is trained on a resample with repeated rows, and their small spread indicates a stable
training procedure. The base paper reported a ROC-AUC of 0.720 for LightGBM on the same dataset~\cite{shreya2025base};
the preprocessing and features differ, so this is an indication rather than a controlled comparison.

\input{tables/models}

Table~\ref{tab:bands} shows how the router divides the test applicants. About 35\% are approved automatically, with an
actual default rate of 2.1\%; 20\% are rejected, with an actual default rate of 21.4\%; the remaining 45\% go to an
officer. In every band the mean predicted risk is within 0.3 percentage points of the actual default rate, so the
probabilities shown to officers and applicants are well calibrated.

\input{tables/bands}

\subsection{Explanations}
Table~\ref{tab:explanations} and Fig.~\ref{fig:explanations} compare the two LIME set-ups with SHAP. LIME applied to
the whole model is effectively unrelated to the applicant: its rank correlation with SHAP is 0.02, it gets the
direction of SHAP's top five features right about half the time, and its locality gap is zero, which means its
explanation agrees as well with another applicant's SHAP values as with the applicant's own. Local LIME raises the
rank correlation to 0.64 and the sign agreement to 73\%, its locality gap of 0.51 shows that its explanations are
specific to the applicant, and the top five features stay the same across seeds in 83\% of cases. The improvement holds
in all three bands (Fig.~\ref{fig:explanations}).

\input{tables/explanations}

\Figure[t!](topskip=0pt, botskip=0pt, midskip=0pt){figures/fig4_explanations.pdf}
{Agreement of whole-model and local LIME with SHAP, by decision band.\label{fig:explanations}}

\subsection{Recourse}
Tables~\ref{tab:recourse_reject_to_review} and~\ref{tab:recourse_review_to_approve} report all methods, and
Fig.~\ref{fig:recourse} summarizes the main ones.

\input{tables/recourse_reject_to_review}
\input{tables/recourse_review_to_approve}

\Figure[t!](topskip=0pt, botskip=0pt, midskip=0pt){figures/fig2_recourse.pdf}
{Share of applicants given a usable recourse option, and one that is also still valid on all five held-out
retrained models, by method and tier.\label{fig:recourse}}

\emph{Coverage and cost.} With the three loan levers, the proposed engine gives a usable option to 57.3\% of
rejected applicants and 37.0\% of Review applicants. Adding the waiting lever raises this to 62.0\% and 45.3\%. Random
search reaches a similar coverage among rejected applicants (59.7\%), but its options change all three levers, whereas
the proposed engine changes 1.5 on average and asks for a smaller change (0.185 against 0.269). The baseline engine,
which has neither the safety margin nor the plausibility check, produces a usable option for only 22.7\% and 12.0\% of
applicants: most of its suggestions either land just past the cutoff or describe an implausible loan.

\emph{Ablations.} Without the safety margin, the engine finds options for most applicants, but only 10.0\% of
rejected applicants receive one that clears the margin, and the share of options that hold on the retrained models
falls from 24.4\% to 13.3\%. Without the plausibility check, 53.3\% receive an option that is both plausible and clears
the margin, against 57.3\% with it, because the cheapest option found is then often implausible. Doubling the margin to
0.2 makes options more durable (34.5\% hold) but reduces coverage to 48.3\%.

\emph{Robustness to retraining.} Standard options are fragile: only 24.4\% of the proposed engine's options for
rejected applicants, and 11.7\% for Review applicants, still reach the promised band on all five held-out retrained
models. Requiring agreement of the five search models raises these shares to 53.3\% and 43.4\%
(Fig.~\ref{fig:robustness}), at the cost of coverage (40.0\% and 17.7\%). The hybrid used by the application keeps the
coverage of the standard search with waiting and gives the largest share of applicants an option that is both usable
and survives retraining: 23.7\% of rejected applicants and 11.0\% of Review applicants, against 17.7\% and 8.7\% for
random search, 14.3\% and 5.0\% for the standard search with waiting, and 8.7\% and 1.7\% for the baseline engine.

\Figure[t!](topskip=0pt, botskip=0pt, midskip=0pt){figures/fig3_robustness.pdf}
{Number of held-out retrained models on which each method's options for rejected applicants still reach the promised
band.\label{fig:robustness}}

\emph{Comparison with DiCE.} On the 20 applicants per tier on which DiCE was run, it produced a usable option for 20\%
of rejected and 5\% of Review applicants, against 60\% and 40\% for the hybrid on the same applicants, and took about
twice as long per applicant (216~ms against 95~ms). Its options also asked for larger changes. These figures rest on
a small sample and are indicative only.

\emph{Sensitivity to the price of waiting.} Qualification does not depend on $\lambda$ in \eqref{eq:cost}, so coverage
is unaffected; only the order of the options can change. Doubling $\lambda$ to 0.08 changes the first option offered for
0.5\% of rejected and 4.4\% of Review applicants, and halving it to 0.02 for 5.9\% and 10.3\%.

\subsection{Age and Fairness}
\label{sec:age}
Automatic approval differs strongly across age bands (Fig.~\ref{fig:age}). On the test set, 13.2\% of applicants
under 25 are approved automatically, against 48.8\% of those over 55, a demographic parity difference of 0.356. For applicants with and without children, the difference is 0.049. The gap
follows a real difference in repayment: 11.4\% of applicants under 25 had payment difficulties, against 5.2\% of those
over 55. The model is calibrated within each age band, with a mean predicted risk within one percentage point of the
actual default rate in every band (Table~\ref{tab:age_experiment}).

\Figure[t!](topskip=0pt, botskip=0pt, midskip=0pt){figures/fig5_age.pdf}
{Actual default rate and automatic approval rate by age band on the test set.\label{fig:age}}

Table~\ref{tab:age_experiment} reports the age experiment. Removing the three direct age features raises the
demographic parity difference from 0.349 to 0.364, and also removing the three features derived from age raises it to
0.377, while the ROC-AUC falls from 0.7668 to 0.7646. Age can be predicted from the remaining features with an $R^2$ of
0.735 (a mean error of 4.5 years), so removing the column does not remove the information. Without the explicit
feature, the model instead overestimates the risk of applicants under 25 (13.8\% predicted against 11.4\% actual) and
approves fewer of them (9.8\% instead of 13.4\%). Equal approval rates are reached only with separate cutoffs for each
age band (difference 0.022), which uses age at decision time and raises the default rate among approved applicants
under 25 from 0.9\% to 3.4\%. This experiment was run before the cleaning statistics were restricted to the training
rows; the rebuilt audit of the current model gives 0.356 instead of 0.349, and the comparison between variants is
unaffected.

\input{tables/age_experiment}

\subsection{System Performance and Understandability}
Table~\ref{tab:latency} and Fig.~\ref{fig:latency} report the run time. An applicant's full path through the
application (submitting, receiving a decision with SHAP, local LIME, recourse and loan terms, reading it and
downloading the PDF report) takes 158~ms on average and 259~ms at the 95th percentile on a laptop. Recourse is the
largest stage (95~ms on average, for Review and Reject cases), followed by local LIME (38~ms) and SHAP (26~ms). Loading
the models at start-up takes 1.6~s.

\input{tables/latency}

\Figure[t!](topskip=0pt, botskip=0pt, midskip=0pt){figures/fig6_latency.pdf}
{Mean time per applicant by stage of the decision.\label{fig:latency}}

Table~\ref{tab:understandability} summarizes the text an applicant reads. It scores 60.0 on the Flesch reading ease
scale and corresponds to a Flesch--Kincaid grade of 7.8, which is plain English. It contains no technical terms or raw
column names, every factor shown has a plain-language name, and an applicant who is not approved sees five reasons and
usually one or two options of about two changes each, in a median of 114 words.

\input{tables/understandability}
